%% ============================================================
%% PRISM Paper 1: supplement
%% Companion to paper.tex
%% ============================================================
\documentclass[10pt,a4paper]{extarticle}
\usepackage{asymmetry}

\newcolumntype{Y}{>{\raggedright\arraybackslash}X}
\setlength{\tabcolsep}{5.2pt}
\renewcommand{\arraystretch}{1.12}

\renewcommand{\thesection}{\Alph{section}}
\renewcommand{\thetable}{S\arabic{table}}
\renewcommand{\thefigure}{S\arabic{figure}}

\newenvironment{claimnotes}
  {\par\vspace{2.5pt}\begin{minipage}{\linewidth}\footnotesize
   {\sffamily\bfseries\color{ACslate}Claim notes.}\enspace
   \setlength{\parskip}{1pt}}
  {\end{minipage}}

\begin{document}
\thispagestyle{firstpage}
\vspace{2pt}

\ACmaketitle%
  {Supplement: PRISM CPU/GPU Portfolio Optimization}%
  {Eligibility Ledger, Scale Frontier, and Workflow Sanity Checks}

\ACauthorline

\noindent
This supplement accompanies \emph{PRISM: A CPU/GPU Portfolio Optimization
Engine for Deadline-Bounded Institutional Rebalancing}. It records material
that supports the main paper without being required for its central claims:
the solver eligibility ledger, the CPU/GPU scale frontier, workflow sanity
checks, and additional artifact details. All evidence-artifact names, timing
lanes, status vocabulary, and claim rules are those defined in the main paper.

\vspace{4pt}

%% ============================================================
\section{Solver Eligibility Ledger}
\label{app:eligibility}
%% ============================================================

Table~\ref{tab:versions} records the full eligibility ledger for the evidence
package, including installed packages that did not produce claim-bearing rows.
Entries are operational eligibility records under the declared protocol, not
solver critiques.

\begin{table}[H]
\centering
\caption{Solver eligibility ledger for the recorded evidence package.}
\label{tab:versions}
\small
\begin{tabularx}{0.97\linewidth}{l l Y}
\toprule
\hdrrow\hcell{Solver stack} & \hcell{Version} & \hcell{Eligibility} \\
\midrule
PRISM-CPU & evidence build \texttt{512758f} & claim-bearing where recorded \\
PRISM-GPU & evidence build \texttt{512758f} & claim-bearing; \texttt{e6} rows are timing/status only \\
OSQP & 1.1.0 & claim-bearing where recorded \\
Clarabel & 0.11.1 & claim-bearing where recorded \\
SCS & via CVXPY 1.8.1 & claim-bearing where recorded \\
MOSEK & 11.1.11 & claim-bearing where recorded \\
CVXPY interface & 1.8.1 & comparator modeling layer; backend named per row \\
Commercial reference (anonymized) & recorded in evidence & objective-gap reference only; timings unpublished under license terms \\
CPLEX & 22.1.2.0 & configured license did not support rows at $N{\geq}500$; completed public-small lane only; not eligible for claim-bearing comparisons \\
HiGHS & 1.13.0 & public-small lane only; not a QP comparator in this package \\
NVIDIA cuOpt & 26.2.0 & no versioned active run recorded; not benchmarked \\
CVXPortfolio & not installed & not benchmarked \\
\bottomrule
\end{tabularx}
\end{table}

%% ============================================================
\section{Scale Frontier Figure and Discussion}
\label{app:frontier}
%% ============================================================

Table~\ref{tab:frontier} and Figure~\ref{fig:gpu} record the full-call
CPU/GPU scale frontier from evidence \texttt{e1\_multisolver\_public\_rows}
and \texttt{e4\_deadline\_completion\_grid}. Full-call timing includes
public-call setup, data movement, status handling, validation, and output
materialization; the frontier is a systems-boundary measurement, not a
device-throughput claim. The convergence of CPU and GPU wall-clock values at
selected sizes is a public-call systems effect, not evidence about
device-level equivalence.

\begin{table}[H]
\centering
\caption{PRISM CPU/GPU scale frontier, full-call wall-clock, evidence
\texttt{e1}/\texttt{e4}. Rows without completed reference solvers carry timing
and feasibility status only.}
\label{tab:frontier}
\small
\begin{tabularx}{0.97\linewidth}{r r r Y}
\toprule
\hdrrow\hcell{$N$} & \hcell{CPU full-call ms} & \hcell{GPU full-call ms} &
\hcell{Reading (systems lane)} \\
\midrule
5{,}000 & 1{,}804.8 & \textbf{140.7} & GPU path is the fastest recorded PRISM path on this row. \\
10{,}000 & 228.3 & \textbf{215.9} & Both PRISM paths complete inside 0.5\,s. \\
25{,}000 & 1{,}063.6 & \textbf{458.3} & GPU path remains inside 1\,s. \\
50{,}000 & 2{,}187.9 & \textbf{861.4} & GPU path remains inside the 5\,s deadline lane. \\
100{,}004 & \textbf{2{,}157.8} & 2{,}161.2 & Public-call systems effect; see text. \\
\bottomrule
\end{tabularx}
\begin{claimnotes}
\texttt{e4} also records GPU memory below 100\,MB at $N{=}100{,}004$, treated
as an externally observed memory-class result.
\end{claimnotes}
\end{table}

\begin{figure}[H]
\centering
\includegraphics[width=0.85\linewidth]{figures/fig_scale_frontier.pdf}
\caption{CPU/GPU scale frontier from Table~\ref{tab:frontier}, evidence
\texttt{e1}/\texttt{e4}. Full-call timing includes public-call setup, data
movement, status handling, validation, and output materialization; the figure
is a systems-boundary measurement, not a device-throughput claim.}
\label{fig:gpu}
\end{figure}

The deadline-completion grid in \texttt{e4} complements the frontier: it
records, for each size and declared deadline from 0.1\,s to 300\,s, whether
each PRISM path completed inside the window. PRISM-GPU completes the
$N{=}10{,}000$ row inside 0.5\,s and the $N{=}100{,}004$ row inside 10\,s in
that grid; both PRISM paths complete every recorded size at the 30\,s and
300\,s deadlines.

%% ============================================================
\section{Walk-Forward and Monte Carlo Sanity Checks}
\label{app:sanity}
%% ============================================================

These checks verify pipeline behavior on repeated portfolio outputs. They are
workflow-correctness evidence under the main paper's claim contract and are
not investment-performance evidence.

\subsection{Walk-Forward Check}

The walk-forward check in
\texttt{e3\_external\_diagnostics\_and\_gap\_checks} covers 131 U.S.
public-equity survivors from January 3, 2019 through December 31, 2025, with
300 weekly rebalance points, 5\,bps transaction costs, and a 5\% position cap.
PRISM-Regime reports Sharpe 0.901 with 95\% interval $[0.770, 1.033]$, Sortino
1.144, and maximum drawdown $-36.4\%$. Equal-Weight reports the same Sharpe in
this check; Min-Vol reports 0.923. Two disclosures bound the interpretation.
First, the universe is survivorship-biased by construction and the
rolling-mean return input is non-predictive; the study is a solver and
workflow correctness check. Second, the tie with Equal-Weight is expected:
under a non-informative return input, converging to broad diversification is
the mathematically correct outcome.

\subsection{Monte Carlo Check}

The Monte Carlo study in \texttt{e3} evaluates 1{,}000 random universes of 20
to 50 stocks. PRISM reports mean Sharpe 5.315 with positive Sharpe in 79.2\%
of trials; Equal-Weight reports 3.974 and 75.1\%; Min-Vol reports 5.095 and
84.0\%. PRISM beats Equal-Weight in 60.4\% of trials and Min-Vol in 54.1\%.
Cross-trial dispersion is large (standard deviation above 6 Sharpe units),
which is why these numbers support workflow stability rather than any
performance ranking.

%% ============================================================
\section{Additional Artifact Details}
\label{app:artifacts}
%% ============================================================

The public artifact repository
(\texttt{AsymmetryComputing/prism-public-evaluation}) uses the following
layout. PRISM implementation code is not part of the public artifact.

\begin{table}[H]
\centering
\caption{Public artifact repository layout.}
\label{tab:repo}
\small
\begin{tabularx}{0.95\linewidth}{l Y}
\toprule
\hdrrow\hcell{Path} & \hcell{Contents} \\
\midrule
\texttt{README.md}, \texttt{MANIFEST.yml}, \texttt{CITATION.cff}, \texttt{LICENSE} & Index, artifact manifest, citation metadata, license. \\
\texttt{environment/} & Hardware/runtime disclosure and package versions. \\
\texttt{data\_sanitized/} & Benchmark input manifest; no client data, no engine state. \\
\texttt{results/} & One sanitized CSV per evidence artifact (\texttt{e1}--\texttt{e6}). \\
\texttt{figures/} & Published figures (vector). \\
\texttt{tables/} & Machine-readable copies of the published tables. \\
\texttt{validation/} & Public residual and feasibility-gate scripts; table reproduction script. \\
\texttt{evidence\_ledger/} & Claim-to-evidence ledger, artifact manifest, artifact policy. \\
\texttt{paper/}, \texttt{supplement/} & arXiv version and sources. \\
\bottomrule
\end{tabularx}
\end{table}

Each claim-bearing number in the main paper maps to a row in
\texttt{evidence\_ledger/claim\_to\_evidence\_ledger.md}, which records claim
ID, claim text, section, table or figure, evidence artifact, allowed claim
type, disallowed interpretation, and reproducibility level.

%% ============================================================
\section{Extra Tables and Claim Notes}
\label{app:extra}
%% ============================================================

\textbf{Small-problem objective agreement
(\texttt{e2\_small\_problem\_agreement}).} On a public FF30 small-problem
study (five seeds, $N{=}30$), all completed comparators agree on the objective
value to approximately seven significant figures, anchoring the
objective-agreement methodology used for gap eligibility. PRISM timing rows
are not recorded in \texttt{e2}; the artifact is comparator-agreement
evidence only.

\textbf{Status vocabulary.} The badges used in the main paper's tables
abbreviate the declared status vocabulary: \emph{feas} (feasible /
non-certified), \emph{t/o} (timeout at the declared ceiling), and \emph{late}
(completed status after the operating window). ``No accepted feasible
return'' records that no output passed the external feasibility gates before
the deadline; it is not a certified infeasibility status.

\textbf{Pending evidence.} A burden-specific PRISM-GPU rerun of the
constraint-burden curve and a multi-seed rerun of the single-shot
\texttt{e1} rows are recorded as planned evidence extensions.

\end{document}
